\documentclass[11pt]{article}
\usepackage[margin=1in]{geometry}
\usepackage{amsmath,amssymb,amsthm,mathtools}
\usepackage{graphicx}
\usepackage{microtype}
\usepackage[colorlinks=true,linkcolor=blue,citecolor=blue,urlcolor=blue]{hyperref}

\newtheorem{theorem}{Theorem}
\newtheorem{lemma}[theorem]{Lemma}
\newtheorem{corollary}[theorem]{Corollary}
\newcommand{\C}{\mathbb C}
\newcommand{\Pfr}{\mathfrak P}
\newcommand{\ran}{\operatorname{ran}}
\newcommand{\Span}{\operatorname{span}}
\newcommand{\id}{\mathrm{id}}

\title{Quiver-Unitary Invariance Characterizes\\Hyper-Reinhardt Free Spectrahedra}
\author{Automated mathematics run \texttt{fa\_banach\_001}, lane 10}
\date{13 August 2026}

\begin{document}
\maketitle

\begin{center}
\fbox{\begin{minipage}{0.87\textwidth}
\textbf{Status.} Candidate full solution, likely valid, pending human review.
The matrix-level invariance in Remark 4.4 of McCullough is proved equivalent
to the defining path-block form of a hyper-Reinhardt pencil.
\end{minipage}}
\end{center}

\section{Source question}

For a tuple $B=(B_1,\ldots,B_g)$ on a finite-dimensional Hilbert space
$H$, put
\[
 L_B(X)=I+\sum_{j=1}^g
 \bigl(B_j\otimes X_j+B_j^*\otimes X_j^*\bigr),
 \qquad
 \Pfr_B[n]=\{X:L_B(X)\succ0\}.
\]
McCullough calls a free spectrahedron hyper-Reinhardt when it has a defining
tuple of path-block form: there is an orthogonal decomposition
$H=H_0\oplus\cdots\oplus H_g$ for which $B_j$ has only its
$(j-1,j)$ block nonzero.  Lemma 4.3 of \cite{McC} shows that such a domain is
invariant under
\begin{equation}\label{action}
 X_j\longmapsto W_{j-1}^*X_jW_j,
 \qquad W_0,\ldots,W_g\in U(n).
\end{equation}
Remark 4.4 asks whether this invariance characterizes the class.

The final $\Pfr_A$ in the source's displayed question is necessarily a local
typo for $\Pfr_B$: $A$ is not attached to the arbitrary domain being tested,
and the preceding lemma makes the intended statement clear.

\clearpage
\begin{figure}[p]
\centering
\includegraphics[width=\textwidth]{figures/open_problem_crop.png}
\caption{Official arXiv:2107.11641v4 PDF page 17: all of Remark 4.4 and the
block-conjugation proof of the forward implication.}
\end{figure}
\clearpage

\section{Proof intuition}

The symmetry is stronger than ordinary coordinatewise rotations because the
$W_j$ are full matrices and the right unitary on $X_j$ is the left unitary on
$X_{j+1}$.  To exploit that strength, fix a matrix size $n$ and replace the
entries of each $X_j$ by new free variables.  The original and transformed
pencils become two monic pencils in $gn^2$ variables.  The hypothesis says
their free spectrahedra agree at \emph{all} auxiliary levels, so the Linear
Gleichstellensatz makes their minimal coefficient tuples unitarily
equivalent.

At $n=2$, the identity $E_{11}E_{21}=0$ can then be transported through this
unitary equivalence.  If two vertex unitaries in \eqref{action} are
independent, choose their quotient to be the flip $F$; now
$E_{11}FE_{21}=E_{11}\ne0$.  The only way the transported product can still
vanish is for its coefficient-space factor to vanish.  Three placements of
adjoints give exactly the orthogonality relations of a path representation.
Their range and source spaces then reconstruct $H_0,\ldots,H_g$.

\section{Complete characterization}

A monic defining tuple is \emph{minimal} if no proper reducing subpencil
defines the same free spectrahedron (equivalently, no smaller monic pencil
does).  Every free spectrahedron has a minimal defining
tuple, unique up to unitary equivalence; we use the Linear Gleichstellensatz
in this standard form \cite[Proposition 2.2]{EHKM}.  The same tool is used in
the graph characterization of Reinhardt pencils in \cite[Theorem 3.1]{MT}.

\begin{theorem}[answer to Remark 4.4]\label{main}
Let $B=(B_1,\ldots,B_g)$ be a minimal monic defining tuple, with every
$B_j\ne0$.  The following are equivalent.
\begin{enumerate}
\item For every $n$, every $X\in\Pfr_B[n]$, and every
$W_0,\ldots,W_g\in U(n)$, the tuple in \eqref{action} belongs to
$\Pfr_B[n]$.
\item There is an orthogonal decomposition
\[
 H=H_0\oplus H_1\oplus\cdots\oplus H_g
\]
such that, with $P_r$ the projection onto $H_r$,
\begin{equation}\label{pathform}
 B_j=P_{j-1}B_jP_j\qquad(1\le j\le g).
\end{equation}
\end{enumerate}
Thus quiver-unitary invariance characterizes hyper-Reinhardt free
spectrahedra.  If the one-coordinate slices have radius one, the nonzero
blocks in \eqref{pathform} have norm one, exactly matching the source's
normalization.
\end{theorem}

We first verify the one technical point needed before applying the
Gleichstellensatz.

\begin{lemma}[matrix-unit inflation preserves minimality]\label{inflation}
Let $B$ be minimal and let $(E_{ab})_{a,b=1}^n$ be the matrix units of $M_n$.
Then the tuple
\[
 C_{j,ab}=B_j\otimes E_{ab}\qquad
 (1\le j\le g,\ 1\le a,b\le n)
\]
is a minimal defining tuple in $gn^2$ variables.
\end{lemma}

\begin{proof}
Decompose $B$ into irreducible reducing summands $B^{(\nu)}$.  Minimality
means that no such summand is redundant.  Fix one summand.  By irreducibility,
the unital $C^*$-algebra generated by its coefficients is the full matrix
algebra $M_{d_\nu}$.  The span of positive-length $*$-words in the
$B_j^{(\nu)}$ is a nonzero two-sided ideal of $M_{d_\nu}$, hence is all of
$M_{d_\nu}$.

Every positive-length $*$-word $Q$ can be lifted, using a compatible chain of
matrix units, to $Q\otimes E_{ab}$ in the $*$-algebra generated by the
$C_{j,ab}^{(\nu)}$.  Therefore that algebra contains
$M_{d_\nu}\otimes M_n$ and the inflated summand is irreducible.

It remains to exclude redundancy among these summands.  Restrict the new
variables by setting $Y_{j,aa}=Z_j$ for every $a$ and
$Y_{j,ab}=0$ for $a\ne b$.  Then
\[
 \sum_{a,b}E_{ab}\otimes Y_{j,ab}=I_n\otimes Z_j,
\]
and, after a tensor-factor permutation, each inflated summand becomes $n$
copies of the original summand.  Redundancy after inflation would therefore
imply redundancy before inflation, a contradiction.  Thus $C$ is minimal.
\end{proof}

\begin{lemma}[forced products]\label{products}
Under condition (1) of Theorem \ref{main},
\begin{align}
 B_jB_k&=0 &&\text{if }k\ne j+1,\label{same}\\
 B_jB_k^*&=0 &&\text{if }j\ne k,\label{right}\\
 B_j^*B_k&=0 &&\text{if }j\ne k.\label{left}
\end{align}
\end{lemma}

\begin{proof}
Fix $n$ and write
\[
 X_j(Y)=\sum_{a,b=1}^n E_{ab}\otimes Y_{j,ab}.
\]
The pencil $L_B(X(Y))$, after permuting tensor factors, has coefficients
\[
 C_{j,ab}=B_j\otimes E_{ab}.
\]
For fixed $W_0,\ldots,W_g\in U(n)$, the transformed pencil has coefficients
\[
 D_{j,ab}=B_j\otimes W_{j-1}^*E_{ab}W_j.
\]
Condition (1), applied at sizes $nm$ for every auxiliary size $m$, says
precisely that $C$ and $D$ define the same free spectrahedron.  Both are
minimal by Lemma \ref{inflation}; hence the Linear Gleichstellensatz gives a
unitary $U_W$ such that
\begin{equation}\label{intertwine}
 U_WC_{j,ab}U_W^*=D_{j,ab}
\end{equation}
for every $j,a,b$.

Now take $n=2$, let $F=\left(\begin{smallmatrix}0&1\\1&0\end{smallmatrix}\right)$,
and note
\begin{equation}\label{witness}
 E_{11}E_{21}=0,
 \qquad E_{11}FE_{21}=E_{11}\ne0.
\end{equation}
If $k\ne j+1$, then $W_j$ and $W_{k-1}$ are independent choices.  We have
$C_{j,11}C_{k,21}=0$; choose their quotient
$W_jW_{k-1}^*=F$.  By \eqref{witness}, the matrix factor of
$D_{j,11}D_{k,21}$ is nonzero (the remaining outer factors are unitary).
Equation \eqref{intertwine} carries the zero product to a zero product, so
$B_jB_k=0$.  This proves \eqref{same}.

For $j\ne k$, start instead with
\[
 C_{j,11}C_{k,12}^*=0
 \quad\text{and choose }W_jW_k^*=F;
\]
the same argument gives \eqref{right}.  Finally, start with
\[
 C_{j,11}^*C_{k,21}=0
 \quad\text{and choose }W_{j-1}W_{k-1}^*=F
\]
to obtain \eqref{left}.  All choices are available exactly in the stated
off-diagonal cases.
\end{proof}

\begin{proof}[Proof of Theorem \ref{main}]
Assume (1).  Put
\[
 R_j=\ran B_j,
 \qquad S_j=\ran B_j^*.
\]
Relation \eqref{left} makes the $R_j$ pairwise orthogonal, and
\eqref{right} makes the $S_j$ pairwise orthogonal.  Moreover,
\eqref{same} gives
\[
 R_k\perp S_j\qquad\text{unless }k=j+1,
\]
because
$\langle B_kx,B_j^*y\rangle=\langle B_jB_kx,y\rangle$.
Consequently the subspaces
\begin{equation}\label{layers}
 H_0=R_1,\qquad
 H_r=S_r+R_{r+1}\ (1\le r\le g-1),\qquad
 H_g=S_g
\end{equation}
are pairwise orthogonal.

Let $K$ be the orthogonal complement of their sum.  Since $K$ is orthogonal
to every $S_j$ and every $R_j$, both $B_j$ and $B_j^*$ vanish on $K$.
Thus $K$ is a common reducing zero summand.  If $K\ne0$, deleting the identity
block on $K$ gives a smaller defining pencil, contrary to minimality.  Hence
the spaces in \eqref{layers} sum to $H$.

In finite dimension, $(\ker B_j)^\perp=S_j$, so $B_j$ vanishes on
$S_j^\perp$ and has range $R_j$.  Since $S_j\subseteq H_j$ and
$R_j\subseteq H_{j-1}$, equation \eqref{pathform} follows.  This proves (2).

Conversely assume (2) and set
\[
 \mathcal W=\bigoplus_{r=0}^g(I_{H_r}\otimes W_r).
\]
A block computation gives
\[
 \mathcal W^*L_B(X)\mathcal W
 =L_B(W_0^*X_1W_1,\ldots,W_{g-1}^*X_gW_g).
\]
Positive definiteness is therefore preserved, proving (1).

For the normalization statement, set all coordinates except $X_j$ equal to
zero.  The only nontrivial block is
\[
 \begin{pmatrix}I&C_j\otimes X_j\\C_j^*\otimes X_j^*&I\end{pmatrix},
 \qquad C_j=B_j|_{H_j}:H_j\to H_{j-1},
\]
which is positive exactly when $\|C_j\|\,\|X_j\|<1$.  A coordinate-slice
radius of one is therefore equivalent to $\|C_j\|=1$.
\end{proof}

\begin{corollary}[inactive coordinates]
If some coordinate is absent from a minimal pencil, the characterization
holds after deleting it, with zero connecting blocks reinserted.  Thus the
unqualified structural statement remains true if zero-dimensional layers or
zero connectors are admitted; the source's norm-one convention simply
excludes these unbounded coordinate directions.
\end{corollary}

\section{Verification, novelty, and limitations}

The bundled exact checker verifies the $2\times2$ witness, the three families
of zero-product relations on a nontrivial three-arrow path pencil, and the
block-conjugation identity.  These checks guard the index and adjoint
placements; they do not replace Lemma \ref{inflation} or the
Gleichstellensatz.

Eight focused attempts are recorded in the run attempt note.  A bounded
literature search on 13 August 2026 covered the exact title, Remark 4.4 phrase,
arXiv id, authors, and combinations of ``hyper-Reinhardt'', ``unitary
invariance'', ``quiver'', and ``characterization''.  No later claimed answer
was found.  McCullough--Tuovila \cite{MT} characterize only the weaker scalar
torus symmetry.  Ben Jemaa--McCullough \cite{BJM} later use \eqref{action} as
a necessary obstruction for a particular domain, but do not state the
converse.  Novelty confidence is moderate because the search was bounded.

Human review should focus on Lemma \ref{inflation} and on the hypotheses of
the Linear Gleichstellensatz.  No conditional mathematical dependency is
currently known.  The only statement caveats are the source's obvious
$A/B$ typo and its harmless norm-one coordinate normalization.

\begin{thebibliography}{9}
\bibitem{McC}
S. McCullough,
\emph{Hyper-Reinhardt Free Spectrahedra},
arXiv:2107.11641v4; published as a chapter in
\emph{Recent Developments in Operator Theory, Mathematical Physics and
Complex Analysis}, 2023.

\bibitem{EHKM}
E. Evert, J. W. Helton, I. Klep, and S. McCullough,
\emph{Circular free spectrahedra},
J. Math. Anal. Appl. 445 (2017), 1047--1070;
arXiv:1604.05756.

\bibitem{MT}
S. McCullough and N. Tuovila,
\emph{Reinhardt free spectrahedra},
Linear Algebra Appl. 640 (2022), 91--117;
arXiv:2012.02289.

\bibitem{BJM}
M. Ben Jemaa and S. McCullough,
\emph{Reinhardt Free Spectrahedra},
arXiv:2301.02746v2 (2023).
\end{thebibliography}

\end{document}
